\documentclass[11pt,a4paper]{article}
\usepackage{estilo-libro}

% ---------- Atajos propios de este documento ----------
% \ar[estilo]{desde}{hasta}{etiqueta}: dibuja un arco con su etiqueta flujo/capacidad.
% Para mover la etiqueta se pasa lbl/.style={pos=...} dentro del estilo.
\tikzset{lbl/.style={}, libre/.style={arco, gray!70}}
\newcommand{\ar}[4][arco]{\draw[#1] (#2) to node[etq, lbl] {#4} (#3);}

% Nodos de cada red (posiciones de la calculadora; el eje y se invierte)
\newcommand{\nodosSunco}{%
  \node[fuente]   (so) at (0,-1)   {so};
  \node[nodo]     (n1) at (1.4,0)  {1};
  \node[nodo]     (n2) at (1.4,-2) {2};
  \node[nodo]     (n3) at (2.8,0)  {3};
  \node[sumidero] (si) at (4.2,-1) {si};}
\newcommand{\nodosSeervada}{%
  \node[fuente]   (O) at (0,-2)     {O};
  \node[nodo]     (A) at (1.5,-0.6) {A};
  \node[nodo]     (B) at (2,-2)     {B};
  \node[nodo]     (C) at (1.5,-3.4) {C};
  \node[nodo]     (D) at (3.6,-0.9) {D};
  \node[nodo]     (E) at (3.6,-3.1) {E};
  \node[sumidero] (T) at (5,-2)     {T};}
\newcommand{\nodosTaha}{%
  \node[fuente]   (t1) at (0,-2)     {1};
  \node[nodo]     (t2) at (1.8,-0.4) {2};
  \node[nodo]     (t3) at (1.8,-2)   {3};
  \node[nodo]     (t4) at (1.8,-3.6) {4};
  \node[sumidero] (t5) at (3.6,-2)   {5};}

\title{\textbf{El problema de flujo máximo}\\[4pt]\large Explicación del método y ejercicios resueltos}
\author{Investigación de Operaciones · Flujo en redes}
\date{}

\begin{document}
\maketitle
\tableofcontents
\newpage

% =====================================================================
\part{¿En qué consiste el problema de flujo máximo?}

\section{¿Qué tipo de problema resuelve?}

Muchas cosas viajan por \textbf{redes}: el agua por las tuberías de una ciudad, los carros por las calles, los
datos por los cables de internet o las personas por los pasillos de un edificio cuando hay que evacuarlo. Cada
tramo de la red tiene un \textbf{límite}: una tubería solo deja pasar cierta cantidad de litros por minuto y una
calle solo deja pasar cierta cantidad de carros por hora.

El problema de flujo máximo pregunta: \textbf{¿cuánto es lo máximo que se puede enviar desde un punto de partida
hasta un punto de llegada, sin pasarse del límite de ningún tramo?}

Para entenderlo usaremos un ejemplo pequeño, que en la Parte II se resuelve completo:

\begin{ojo}
Una petrolera manda petróleo desde un pozo (\textbf{so}) hasta una refinería (\textbf{si}) por una red de
oleoductos que pasa por tres estaciones de bombeo (1, 2 y 3). Cada tubería aguanta un máximo, en millones de
barriles por hora: so$\to$1 aguanta 2, so$\to$2 aguanta 3, 1$\to$2 aguanta 3, 1$\to$3 aguanta 4,
3$\to$si aguanta 1 y 2$\to$si aguanta 2. ¿Cuánto petróleo puede llegar a la refinería por hora como máximo?
\end{ojo}

Todo problema de flujo máximo tiene estas partes:

\begin{center}\small
\begin{tabular}{p{3.0cm}p{6.6cm}p{4.6cm}}
\toprule
\textbf{Parte} & \textbf{Qué es} & \textbf{En el ejemplo} \\
\midrule
Nodos & Los puntos de la red (cruces, estaciones, ciudades). & so, 1, 2, 3, si \\
Arcos & Los tramos que unen dos nodos, con un sentido (la flecha). & las 6 tuberías \\
Capacidad & Lo máximo que puede pasar por un arco. & so$\to$1 aguanta 2, etc. \\
Fuente & El nodo de donde sale todo lo que se envía. & el pozo so \\
Sumidero & El nodo a donde tiene que llegar. & la refinería si \\
Flujo & Lo que de verdad se manda por cada arco. & lo que hay que decidir \\
\bottomrule
\end{tabular}
\end{center}

Además, el flujo tiene que cumplir \textbf{dos reglas}:
\begin{enumerate}[itemsep=2pt]
  \item \textbf{No pasarse de la capacidad:} por cada arco pasa entre 0 y su capacidad.
  \item \textbf{Conservación:} en cada nodo intermedio, \textbf{todo lo que entra, sale}. Las estaciones de
        bombeo no guardan ni fabrican petróleo; solo lo dejan pasar.
\end{enumerate}

\begin{center}
\begin{tikzpicture}[x=2.4cm, y=1.6cm]
  \nodosSunco
  \ar{so}{n1}{2}
  \ar{so}{n2}{3}
  \ar{n1}{n2}{3}
  \ar{n1}{n3}{4}
  \ar{n3}{si}{1}
  \ar{n2}{si}{2}
  \node[font=\scriptsize, below=1mm of so] {fuente};
  \node[font=\scriptsize, below=1mm of si] {sumidero};
\end{tikzpicture}

\small La red del ejemplo. El número de cada arco es su capacidad.
\end{center}

% =====================================================================
\section{La idea del método}

Imaginemos un tubo formado por pedazos de distinto grosor. Por más agua que se meta por un extremo, por el otro
solo sale lo que deja pasar \textbf{el pedazo más angosto}. A ese pedazo se le llama \textbf{cuello de botella}.

\begin{center}
\begin{tikzpicture}[x=1cm, y=1cm]
  \fill[gray!25] (0,-0.9) rectangle (3,0.9);
  \fill[gray!25] (3,-0.3) rectangle (5,0.3);
  \fill[gray!25] (5,-0.6) rectangle (8,0.6);
  \draw[thick] (0,0.9)--(3,0.9)--(3,0.3)--(5,0.3)--(5,0.6)--(8,0.6);
  \draw[thick] (0,-0.9)--(3,-0.9)--(3,-0.3)--(5,-0.3)--(5,-0.6)--(8,-0.6);
  \node[font=\small] at (1.5,0) {cabe 3};
  \node[font=\small\bfseries] at (4,0) {cabe 1};
  \node[font=\small] at (6.5,0) {cabe 2};
  \draw[-{Stealth}, thick] (-1.2,0)--(-0.1,0) node[pos=0, left, font=\small] {entra};
  \draw[-{Stealth}, thick] (8.1,0)--(9.2,0) node[right, font=\small] {sale 1};
  \draw[-{Stealth}] (4,-1.4) node[below, font=\scriptsize] {cuello de botella} -- (4,-0.4);
\end{tikzpicture}
\end{center}

Una red es como muchos tubos de ese tipo, unidos unos con otros. El método más conocido para resolver el problema
es el de \textbf{Ford y Fulkerson} (1956), y su idea es muy sencilla:

\begin{enumerate}[itemsep=3pt]
  \item Empezar sin mandar nada: todo el flujo vale 0.
  \item Buscar un \textbf{camino} de la fuente al sumidero por el que \textbf{todavía quepa algo}.
  \item Mandar por ese camino todo lo que permite su tramo más angosto (el cuello de botella).
  \item Repetir hasta que ya no exista ningún camino con espacio. Lo que se mandó en total es el flujo máximo.
\end{enumerate}

En 1972, \textbf{Edmonds y Karp} agregaron una regla para elegir el camino: \textbf{usar siempre el camino con
menos arcos}. Para encontrarlo se revisa la red ``por capas'': primero los nodos a un paso de la fuente, luego los
que están a dos pasos, y así sucesivamente (a esto se le llama \textbf{búsqueda en anchura}). Con esta regla el
método siempre termina, y el número de caminos que hay que buscar depende solo del tamaño de la red, no de lo
grandes que sean las capacidades. Es la versión que usa la calculadora y la que se sigue en este documento.

% =====================================================================
\section{¿Por qué a veces hay que devolver flujo?}

Un camino mal escogido puede ``tapar'' la red. Veamos una red de cuatro nodos donde todos los arcos aguantan 1.
Si primero se manda 1 unidad por $s\to a\to b\to t$, ya no queda ningún camino ``hacia adelante'': $s\to a$ y
$b\to t$ están llenos. Parece que el máximo es 1, pero no es cierto: se pueden mandar 2.

\begin{center}
\begin{tikzpicture}[x=1.9cm, y=1.0cm]
  \node[fuente] (s) at (0,0) {$s$};
  \node[nodo] (a) at (1.5,1) {$a$};
  \node[nodo] (b) at (1.5,-1) {$b$};
  \node[sumidero] (t) at (3,0) {$t$};
  \ar[camino]{s}{a}{1/1}
  \ar[libre]{s}{b}{0/1}
  \ar[camino]{a}{b}{1/1}
  \ar[libre]{a}{t}{0/1}
  \ar[camino]{b}{t}{1/1}
  \node[font=\small, align=center] at (1.5,-2) {Primer camino: $s\to a\to b\to t$\\ (flujo total 1)};
\end{tikzpicture}
\qquad
\begin{tikzpicture}[x=1.9cm, y=1.0cm]
  \node[fuente] (s) at (0,0) {$s$};
  \node[nodo] (a) at (1.5,1) {$a$};
  \node[nodo] (b) at (1.5,-1) {$b$};
  \node[sumidero] (t) at (3,0) {$t$};
  \ar[usado]{s}{a}{1/1}
  \ar[camino]{s}{b}{1/1}
  \ar[libre]{a}{b}{0/1}
  \draw[camino, dashed] (b) to[bend left=40] node[etq, left] {devuelvo 1} (a);
  \ar[camino]{a}{t}{1/1}
  \ar[usado]{b}{t}{1/1}
  \node[font=\small, align=center] at (1.5,-2) {Segundo camino: $s\to b\to a\to t$\\ (flujo total 2)};
\end{tikzpicture}
\end{center}

El truco es permitir que un camino recorra un arco \textbf{al revés}, cuando ese arco ya lleva flujo. Recorrer
$a\to b$ al revés significa \textbf{devolver} lo que se había mandado por él: la unidad que llegaba a $a$ ya no
sigue hacia $b$, sino que se va directo a $t$; y la unidad nueva que llega a $b$ desde $s$ ocupa su lugar en
$b\to t$. Al final $a\to b$ queda en 0 y llegan 2 unidades a $t$.

Para tener en cuenta estas devoluciones, el método no busca los caminos en la red original sino en la
\textbf{red residual}, que en cada arco anota dos cosas:
\begin{itemize}[itemsep=2pt]
  \item \textbf{Lo que todavía cabe} hacia adelante: capacidad $-$ flujo.
  \item \textbf{Lo que se puede devolver} hacia atrás: el flujo que ya lleva.
\end{itemize}

% =====================================================================
\section{Las palabras que se usan}

\begin{center}\small
\begin{tabular}{p{3.4cm}p{10.8cm}}
\toprule
\textbf{Término} & \textbf{Qué significa} \\
\midrule
Red & Un conjunto de nodos unidos por arcos con sentido. \\
Capacidad $u_{ij}$ & Lo máximo que puede pasar por el arco que va del nodo $i$ al nodo $j$. \\
Flujo $f_{ij}$ & Lo que se manda por ese arco. Siempre $0 \le f_{ij} \le u_{ij}$. \\
Fuente y sumidero & El nodo de salida y el nodo de llegada. \\
Conservación & En cada nodo intermedio, lo que entra es igual a lo que sale. \\
Red residual & La red de ``lo que todavía cabe'' (hacia adelante) y ``lo que se puede devolver'' (hacia atrás). \\
Camino de aumento & Un camino de la fuente al sumidero, dentro de la red residual, donde cabe algo en cada tramo. \\
Cuello de botella & El tramo del camino con menos espacio. Es lo que se puede mandar por ese camino. \\
Corte & Una línea que separa la red en dos lados: uno con la fuente y otro con el sumidero. \\
Capacidad del corte & La suma de las capacidades de los arcos que cruzan el corte \textbf{del lado de la fuente al lado del sumidero}. \\
Corte mínimo & El corte de menor capacidad: el ``embudo más angosto'' de toda la red. \\
\bottomrule
\end{tabular}
\end{center}

% =====================================================================
\section{Los pasos del método}\label{sec:algoritmo}

Cada vuelta del método (cada \textbf{iteración}) hace cuatro pasos, siempre en el mismo orden. Para que se
entiendan, se muestra lo que ocurre en la primera iteración del ejemplo de la petrolera.

\paso{Paso 0. Empezar sin mandar nada}
Todos los flujos valen 0. Así se cumplen las dos reglas: nada se pasa de su capacidad y en cada nodo entra 0 y sale 0.

\textit{En el ejemplo:} flujo total $= 0$; en la red residual, en cada tubería cabe toda su capacidad.

\paso{Paso 1. Buscar el camino con menos arcos}
En la red residual se busca, por capas, un camino de la fuente al sumidero donde \textbf{cada tramo tenga espacio}
(algo que todavía cabe hacia adelante, o algo que se puede devolver hacia atrás).
\begin{itemize}[itemsep=2pt]
  \item Si \textbf{no hay ningún camino}, el método termina.
  \item Si hay, se toma el de menos arcos. Si hay empate, gana el que se encuentra primero revisando los arcos en
        el orden de la lista.
\end{itemize}
\textit{En el ejemplo:} a un paso de so están 1 y 2; a dos pasos, 3 (desde 1) y si (desde 2). El camino más corto
es so$\to$2$\to$si.

\paso{Paso 2. Calcular el cuello de botella}
Se anota cuánto cabe en cada tramo del camino y se toma el menor.

\textit{En el ejemplo:} en so$\to$2 caben 3 y en 2$\to$si caben 2, así que el cuello de botella es
$\min(3, 2) = 2$.

\paso{Paso 3. Mandar el flujo y actualizar la red residual}
Por cada tramo del camino se suma el cuello de botella al flujo (o se resta, si el tramo era una devolución).
El flujo total aumenta en esa cantidad. Después se vuelve al paso 1.

\textit{En el ejemplo:} so$\to$2 queda en 2 de 3 y 2$\to$si en 2 de 2 (lleno). El flujo total pasa a ser 2.

\begin{idea}[Resumen de los pasos]
\begin{enumerate}[leftmargin=*, itemsep=3pt]
  \item[\textbf{0.}] Empezar con todos los flujos en 0.
  \item \textbf{Buscar camino:} en la red residual, el de menos arcos de la fuente al sumidero.
        Si no hay, terminar.
  \item \textbf{Cuello de botella:} $\delta = $ lo menor que cabe en los tramos del camino.
  \item \textbf{Aumentar:} sumar $\delta$ en los tramos hacia adelante, restarlo en las devoluciones y
        sumar $\delta$ al flujo total. Volver al paso 1.
\end{enumerate}
\end{idea}

\begin{center}
\begin{tikzpicture}[node distance=4mm and 9mm, >=Stealth,
  caja/.style={draw, rounded corners=2pt, fill=gray!8, align=center, minimum width=5.6cm, inner sep=5pt, font=\small},
  decision/.style={draw, diamond, aspect=2.4, fill=gray!18, align=center, inner sep=1pt, font=\small},
  fin/.style={draw, rounded corners=8pt, fill=gray!25, align=center, minimum width=2.8cm, inner sep=6pt, font=\small}]
  \node[caja] (ini) {Todos los flujos en 0};
  \node[caja, below=of ini] (res) {Armar la red residual\\ (lo que cabe y lo que se puede devolver)};
  \node[decision, below=of res] (hay) {¿Hay camino de la\\ fuente al sumidero?};
  \node[fin, right=of hay] (fin) {\textbf{Flujo máximo}\\ y corte mínimo};
  \node[caja, below=of hay] (cue) {Tomar el de menos arcos\\ y calcular su cuello de botella};
  \node[caja, below=of cue] (aum) {Mandar el cuello de botella\\ por el camino};
  \draw[->] (ini) -- (res);
  \draw[->] (res) -- (hay);
  \draw[->] (hay) -- node[above, font=\scriptsize] {no} (fin);
  \draw[->] (hay) -- node[right, font=\scriptsize] {sí} (cue);
  \draw[->] (cue) -- (aum);
  \draw[->] (aum.west) -- ++(-0.9,0) |- node[pos=0.25, left, font=\scriptsize, align=center] {siguiente\\iteración} (res.west);
\end{tikzpicture}
\end{center}

% =====================================================================
\section{¿Cómo se sabe que terminó?}

\begin{description}[leftmargin=0pt, itemsep=5pt]
  \item[Ya no hay camino.] En el paso 1, la búsqueda por capas no llega al sumidero. Desde ese momento
        ningún camino puede llevar nada más, así que el flujo total es el máximo.
  \item[El corte mínimo sale gratis.] Los nodos a los que \textbf{todavía se llega} desde la fuente forman un lado
        de la red; los demás, el otro. Todos los arcos que van del lado de la fuente al otro lado están llenos (y los
        que vuelven en sentido contrario van vacíos), y la suma de sus capacidades es igual al flujo máximo.
  \item[Red sin camino desde el inicio.] Si ni siquiera con todo en 0 hay camino de la fuente al sumidero, el
        flujo máximo es 0.
  \item[Varias soluciones.] El \emph{valor} del flujo máximo es único, pero a veces hay más de una forma de
        repartirlo entre los arcos.
\end{description}

% =====================================================================
\section{El corte mínimo: el embudo más angosto}

Si se dibuja una línea que separe la fuente del sumidero, \textbf{todo} lo que llega al sumidero tiene que cruzar
esa línea por alguno de los arcos que la atraviesan. Por eso ningún flujo puede ser mayor que la capacidad de
\emph{cualquier} corte.

El \textbf{teorema de flujo máximo y corte mínimo} (Ford y Fulkerson, 1956) dice algo más fuerte:
\[
\text{flujo máximo} = \text{capacidad del corte mínimo}.
\]
Es decir, lo máximo que se puede mandar es exactamente lo que deja pasar el embudo más angosto de la red.

\textit{En el ejemplo:} la línea que deja a si sola de un lado corta los arcos 3$\to$si (1) y 2$\to$si (2).
Su capacidad es $1 + 2 = 3$, y como el método logra mandar 3, ya se sabe que no se puede más.

Esto sirve de \textbf{comprobación}: si el flujo encontrado es igual a la capacidad de algún corte, ese flujo es
el máximo. Además dice \textbf{qué tramos conviene ampliar}: agrandar un arco que no cruza el corte mínimo no
aumenta el flujo máximo.

% =====================================================================
\section{Relación con la programación lineal}

El flujo máximo también se puede escribir como un problema de programación lineal y resolverse con el Simplex
Revisado del primer previo. Las variables son los flujos $f_{ij}$; se maximiza lo que sale de la fuente, cada
flujo va entre 0 y su capacidad ($0 \le f_{ij} \le u_{ij}$) y en cada nodo intermedio lo que entra menos lo que
sale es 0. \textit{En el ejemplo:}
\[
\begin{aligned}
\text{Max } Z &= f_{so,1} + f_{so,2}\\
\text{s.a. } \quad f_{so,1} \le 2,\quad f_{so,2} \le 3,\quad f_{12} &\le 3,\quad f_{13} \le 4,\quad f_{3,si} \le 1,\quad f_{2,si} \le 2\\
f_{so,1} - f_{12} - f_{13} &= 0 \nota{nodo 1}\\
f_{so,2} + f_{12} - f_{2,si} &= 0 \nota{nodo 2}\\
f_{13} - f_{3,si} &= 0 \nota{nodo 3}
\end{aligned}
\]
El Simplex da el mismo valor, pero Ford-Fulkerson aprovecha la forma de red: no necesita matrices, se hace a mano
sobre el dibujo y, si las capacidades son enteras, el flujo de cada arco también sale entero.

% =====================================================================
\section{Ventajas y desventajas}

\begin{center}\small
\begin{tabular}{p{7.2cm}p{7.2cm}}
\toprule
\textbf{Ventajas} & \textbf{Desventajas} \\
\midrule
\begin{itemize}[leftmargin=*, itemsep=2pt, topsep=2pt]
  \item Se entiende con el dibujo: cada paso es ``buscar un camino y llenarlo''.
  \item Con capacidades enteras, el flujo en cada arco sale entero.
  \item Al terminar entrega también el corte mínimo, que sirve de comprobación y muestra dónde ampliar.
  \item Con la regla de Edmonds y Karp siempre termina, y rápido.
\end{itemize}
&
\begin{itemize}[leftmargin=*, itemsep=2pt, topsep=2pt]
  \item Hay que llevar con cuidado la red residual, sobre todo las devoluciones.
  \item En redes grandes hay muchos caminos y hacerlo a mano es largo.
  \item Solo responde ``cuánto cabe'': no tiene en cuenta costos (para eso está el flujo a costo mínimo).
\end{itemize}
\\
\bottomrule
\end{tabular}
\end{center}

% =====================================================================
\newpage
\part{Ejercicios resueltos}
\setcounter{section}{0}

Los tres ejercicios siguientes aplican, uno por uno, los pasos de la sección~\ref{sec:algoritmo}. En cada
iteración se busca el camino con menos arcos, se calcula su cuello de botella y se actualizan los flujos. En los
dibujos, cada arco lleva la etiqueta \textbf{flujo/capacidad}, el camino de la iteración va con línea gruesa y el
corte mínimo con línea discontinua.

% =====================================================================
\section{Sunco Oil: el oleoducto}
\textit{Fuente: Winston, Investigación de operaciones: aplicaciones y algoritmos, capítulo de modelos de redes.}

\subsection*{Planteamiento del problema}
Sunco Oil quiere mandar la mayor cantidad posible de petróleo por hora desde el pozo \textbf{so} hasta la
refinería \textbf{si}. En el camino, el petróleo pasa por algunas de las tres estaciones de bombeo (1, 2 y 3).
Cada tubería tiene un máximo, en millones de barriles por hora:

\begin{center}\small
\begin{tabular}{lcccccc}
\toprule
\textbf{Tubería} & so$\to$1 & so$\to$2 & 1$\to$2 & 1$\to$3 & 3$\to$si & 2$\to$si \\
\midrule
\textbf{Capacidad} & 2 & 3 & 3 & 4 & 1 & 2 \\
\bottomrule
\end{tabular}

\medskip
\begin{tikzpicture}[x=2.0cm, y=1.05cm]
  \nodosSunco
  \ar{so}{n1}{0/2}
  \ar{so}{n2}{0/3}
  \ar{n1}{n2}{0/3}
  \ar{n1}{n3}{0/4}
  \ar{n3}{si}{0/1}
  \ar{n2}{si}{0/2}
\end{tikzpicture}

Red inicial: todo el flujo en 0.
\end{center}

La pregunta: \textbf{¿cuántos millones de barriles por hora pueden llegar como máximo a la refinería?}

\subsection*{Formulación del modelo}
Sea $f_{ij}$ los millones de barriles por hora que se mandan por la tubería de $i$ a $j$. Se busca que lo que sale
del pozo, $Z = f_{so,1} + f_{so,2}$, sea lo más grande posible; cada $f_{ij}$ va entre 0 y su capacidad, y en las
estaciones 1, 2 y 3 lo que entra es igual a lo que sale (el modelo de PL completo está en la Parte I).

\subsection*{Caminos de aumento}
\paso{Camino 1: \ so $\to$ 2 $\to$ si}
Por capas: desde so se llega a 1 y a 2; desde 1 se llega a 3 y desde 2 se llega a si. El camino más corto
tiene 2 arcos. Debajo de cada número se indica el tramo y cuánto cabe en él:
\[
\delta = \min(\underbrace{3}_{\text{so}\to2},\ \underbrace{2}_{2\to\text{si}}) = 2,
\qquad \text{flujo total} = 0 + 2 = 2 \nota{2$\to$si se llena}
\]

\paso{Camino 2: \ so $\to$ 1 $\to$ 3 $\to$ si}
Como 2$\to$si ya está lleno, a si solo se llega desde 3. El camino más corto ahora tiene 3 arcos:
\[
\delta = \min(\underbrace{2}_{\text{so}\to1},\ \underbrace{4}_{1\to3},\ \underbrace{1}_{3\to\text{si}}) = 1,
\qquad \text{flujo total} = 2 + 1 = 3 \nota{3$\to$si se llena}
\]

\begin{center}
\begin{tikzpicture}[x=1.6cm, y=0.95cm]
  \nodosSunco
  \ar[libre]{so}{n1}{0/2}
  \ar[camino]{so}{n2}{2/3}
  \ar[libre]{n1}{n2}{0/3}
  \ar[libre]{n1}{n3}{0/4}
  \ar[libre]{n3}{si}{0/1}
  \ar[camino]{n2}{si}{2/2}
  \node[font=\small] at (2.1,-2.75) {Camino 1 (total 2)};
\end{tikzpicture}
\qquad
\begin{tikzpicture}[x=1.6cm, y=0.95cm]
  \nodosSunco
  \ar[camino]{so}{n1}{1/2}
  \ar[usado]{so}{n2}{2/3}
  \ar[libre]{n1}{n2}{0/3}
  \ar[camino]{n1}{n3}{1/4}
  \ar[camino]{n3}{si}{1/1}
  \ar[usado]{n2}{si}{2/2}
  \node[font=\small] at (2.1,-2.75) {Camino 2 (total 3)};
\end{tikzpicture}
\end{center}

\paso{¿Hay otro camino?}
Desde so todavía se llega a 1 (cabe 1 más) y a 2 (cabe 1 más); desde 1 se llega a 2 y a 3. Pero las dos
tuberías que entran a si van llenas: \textbf{ya no se llega a si}. El método termina.

\subsection*{Resultado}
\[
\res{\text{Flujo máximo} = 3}
\]
\begin{itemize}[itemsep=2pt]
  \item A la refinería llegan como máximo \textbf{3 millones de barriles por hora}.
  \item Flujo por tubería: so$\to$1 = 1, so$\to$2 = 2, 1$\to$2 = 0 (sin usar), 1$\to$3 = 1, 3$\to$si = 1 (llena),
        2$\to$si = 2 (llena).
\end{itemize}

\subsection*{Comprobación}
\begin{itemize}[itemsep=2pt]
  \item \textbf{Capacidades:} $1 \le 2$, \ $2 \le 3$, \ $0 \le 3$, \ $1 \le 4$, \ $1 \le 1$, \ $2 \le 2$ \checkmark
  \item \textbf{Conservación:} nodo 1: entra 1, sale $0 + 1 = 1$ \checkmark; \ nodo 2: entra $2 + 0 = 2$, sale 2
        \checkmark; \ nodo 3: entra 1, sale 1 \checkmark
  \item \textbf{Sale de so} $1 + 2 = 3$ y \textbf{llega a si} $1 + 2 = 3$ \checkmark
  \item \textbf{Corte mínimo:} desde so se alcanzan $\{\text{so}, 1, 2, 3\}$. Los arcos que cruzan son
        3$\to$si (1) y 2$\to$si (2): capacidad $1 + 2 = 3 =$ flujo máximo \checkmark
\end{itemize}

\subsection*{Interpretación gráfica}
\begin{center}
\begin{tikzpicture}[x=2.0cm, y=1.05cm]
  \nodosSunco
  \ar[usado]{so}{n1}{1/2}
  \ar[usado]{so}{n2}{2/3}
  \ar[libre]{n1}{n2}{0/3}
  \ar[usado]{n1}{n3}{1/4}
  \ar[usado, lbl/.style={pos=0.3}]{n3}{si}{1/1}
  \ar[usado]{n2}{si}{2/2}
  \draw[corte] (3.6,0.45) -- (3.6,-2.45) node[below, font=\small] {corte mínimo $= 3$};
  \node[font=\scriptsize] at (2.45,-1.0) {lado de la fuente};
  \node[font=\scriptsize] at (4.3,-1.75) {lado del sumidero};
\end{tikzpicture}
\end{center}
La línea discontinua es el embudo de la red: solo dos tuberías llegan a la refinería y entre las dos aguantan 3.
Por más que se amplíen las demás, no puede llegar más petróleo; para aumentar el envío habría que ampliar
3$\to$si o 2$\to$si. La tubería 1$\to$2 no hace falta con este reparto.

% =====================================================================
\newpage
\section{Parque Seervada: los tranvías}
\textit{Fuente: Hillier y Lieberman, Introducción a la Investigación de Operaciones, capítulo de optimización de redes.}

\subsection*{Planteamiento del problema}
En temporada alta, el parque Seervada lleva a los visitantes en tranvía desde la entrada \textbf{O} hasta el
mirador \textbf{T}, pasando por las estaciones A, B, C, D y E. Para cuidar el parque, cada camino admite un número
máximo de viajes de tranvía por día:

\begin{center}\small
\setlength{\tabcolsep}{4.5pt}
\begin{tabular}{lcccccccccccc}
\toprule
\textbf{Camino} & O$\to$A & O$\to$B & O$\to$C & A$\to$B & A$\to$D & B$\to$C & B$\to$D & B$\to$E & C$\to$E & D$\to$T & E$\to$D & E$\to$T \\
\midrule
\textbf{Viajes} & 5 & 7 & 4 & 1 & 3 & 2 & 4 & 5 & 4 & 9 & 1 & 6 \\
\bottomrule
\end{tabular}

\medskip
\begin{tikzpicture}[x=1.8cm, y=1.12cm]
  \nodosSeervada
  \ar{O}{A}{0/5}
  \ar{O}{B}{0/7}
  \ar{O}{C}{0/4}
  \ar{A}{B}{0/1}
  \ar{A}{D}{0/3}
  \ar{B}{C}{0/2}
  \ar[arco, lbl/.style={pos=0.6}]{B}{D}{0/4}
  \ar[arco, lbl/.style={pos=0.6}]{B}{E}{0/5}
  \ar{C}{E}{0/4}
  \ar{D}{T}{0/9}
  \ar{E}{D}{0/1}
  \ar{E}{T}{0/6}
\end{tikzpicture}

Red inicial: todo el flujo en 0.
\end{center}

La pregunta: \textbf{¿cuántos viajes de tranvía por día pueden llegar como máximo de O a T?}

\subsection*{Formulación del modelo}
$f_{ij}$ = viajes por día del tranvía en el camino de $i$ a $j$.
\[
\begin{aligned}
\text{Max } Z &= f_{OA} + f_{OB} + f_{OC} \nota{lo que sale de la entrada}\\
\text{s.a. } \quad f_{ij} &\le u_{ij} \ \text{(la tabla)}, \qquad f_{ij} \ge 0\\
f_{OA} - f_{AB} - f_{AD} &= 0 \nota{estación A}\\
f_{OB} + f_{AB} - f_{BC} - f_{BD} - f_{BE} &= 0 \nota{estación B}\\
f_{OC} + f_{BC} - f_{CE} &= 0 \nota{estación C}\\
f_{AD} + f_{BD} + f_{ED} - f_{DT} &= 0 \nota{estación D}\\
f_{BE} + f_{CE} - f_{ED} - f_{ET} &= 0 \nota{estación E}
\end{aligned}
\]

\subsection*{Caminos de aumento}
Debajo de cada número se indica el tramo; el número es lo que todavía cabe en él.

\paso{Camino 1: \ $O \to A \to D \to T$}
Desde O se llega a A, B y C; desde A, a D; desde B, a E; desde D, a T. Hay varios caminos de 3 arcos y el primero
que se completa es el que pasa por A.
\[
\delta = \min(\underbrace{5}_{OA},\ \underbrace{3}_{AD},\ \underbrace{9}_{DT}) = 3,
\qquad \text{flujo total} = 0 + 3 = 3 \nota{A$\to$D se llena}
\]

\paso{Camino 2: \ $O \to B \to D \to T$}
A$\to$D está lleno, así que ahora a D se llega desde B.
\[
\delta = \min(\underbrace{7}_{OB},\ \underbrace{4}_{BD},\ \underbrace{6}_{DT}) = 4,
\qquad \text{flujo total} = 3 + 4 = 7 \nota{a D$\to$T le quedaban $9 - 3 = 6$}
\]

\paso{Camino 3: \ $O \to B \to E \to T$}
Los dos arcos que llegan a D desde A y B están llenos. Desde B se llega a E, y desde E a T.
\[
\delta = \min(\underbrace{3}_{OB},\ \underbrace{5}_{BE},\ \underbrace{6}_{ET}) = 3,
\qquad \text{flujo total} = 7 + 3 = 10 \nota{O$\to$B se llena}
\]

\paso{Camino 4: \ $O \to C \to E \to T$}
O$\to$B está lleno; desde C se llega a E y desde E a T.
\[
\delta = \min(\underbrace{4}_{OC},\ \underbrace{4}_{CE},\ \underbrace{3}_{ET}) = 3,
\qquad \text{flujo total} = 10 + 3 = 13 \nota{a E$\to$T le quedaban $6 - 3 = 3$}
\]

\paso{Camino 5: \ $O \to C \to E \to D \to T$}
E$\to$T está lleno, pero desde E todavía se puede subir a D (E$\to$D aguanta 1) y a D$\to$T le queda espacio.
Ya no hay caminos de 3 arcos: el más corto tiene 4.
\[
\delta = \min(\underbrace{1}_{OC},\ \underbrace{1}_{CE},\ \underbrace{1}_{ED},\ \underbrace{2}_{DT}) = 1,
\qquad \text{flujo total} = 13 + 1 = 14 \nota{O$\to$C, C$\to$E y E$\to$D se llenan}
\]

\begin{center}
\begin{tikzpicture}[x=1.3cm, y=0.95cm]
  \nodosSeervada
  \ar[usado]{O}{A}{3/5}
  \ar[camino]{O}{B}{4/7}
  \ar[libre]{O}{C}{0/4}
  \ar[libre]{A}{B}{0/1}
  \ar[usado]{A}{D}{3/3}
  \ar[libre]{B}{C}{0/2}
  \ar[camino, lbl/.style={pos=0.6}]{B}{D}{4/4}
  \ar[libre, lbl/.style={pos=0.6}]{B}{E}{0/5}
  \ar[libre]{C}{E}{0/4}
  \ar[camino]{D}{T}{7/9}
  \ar[libre]{E}{D}{0/1}
  \ar[libre]{E}{T}{0/6}
  \node[font=\small] at (2.5,-4.15) {Después del camino 2 (total 7)};
\end{tikzpicture}
\hfill
\begin{tikzpicture}[x=1.3cm, y=0.95cm]
  \nodosSeervada
  \ar[usado]{O}{A}{3/5}
  \ar[usado]{O}{B}{7/7}
  \ar[camino]{O}{C}{4/4}
  \ar[libre]{A}{B}{0/1}
  \ar[usado]{A}{D}{3/3}
  \ar[libre]{B}{C}{0/2}
  \ar[usado, lbl/.style={pos=0.6}]{B}{D}{4/4}
  \ar[usado, lbl/.style={pos=0.6}]{B}{E}{3/5}
  \ar[camino]{C}{E}{4/4}
  \ar[camino]{D}{T}{8/9}
  \ar[camino]{E}{D}{1/1}
  \ar[usado]{E}{T}{6/6}
  \node[font=\small] at (2.5,-4.15) {Después del camino 5 (total 14)};
\end{tikzpicture}
\end{center}

\paso{¿Hay otro camino?}
Desde O todavía se llega a A (caben 2 más); desde A a B (A$\to$B aguanta 1); desde B a C y a E. Pero todos los
arcos que salen de ese grupo hacia D o T están llenos: A$\to$D, B$\to$D, E$\to$D y E$\to$T. \textbf{Ya no se llega
a T}. El método termina.

\subsection*{Resultado}
\[
\res{\text{Flujo máximo} = 14}
\]
\begin{center}\small
\setlength{\tabcolsep}{4.5pt}
\begin{tabular}{lcccccccccccc}
\toprule
\textbf{Camino} & O$\to$A & O$\to$B & O$\to$C & A$\to$B & A$\to$D & B$\to$C & B$\to$D & B$\to$E & C$\to$E & D$\to$T & E$\to$D & E$\to$T \\
\midrule
\textbf{Flujo} & 3 & 7 & 4 & 0 & 3 & 0 & 4 & 3 & 4 & 8 & 1 & 6 \\
\textbf{Capacidad} & 5 & 7 & 4 & 1 & 3 & 2 & 4 & 5 & 4 & 9 & 1 & 6 \\
\bottomrule
\end{tabular}
\end{center}
Se pueden programar como máximo \textbf{14 viajes de tranvía por día} de la entrada al mirador.

\subsection*{Comprobación}
\begin{itemize}[itemsep=2pt]
  \item \textbf{Capacidades:} en la tabla, cada flujo es menor o igual que su capacidad \checkmark
  \item \textbf{Conservación:}
        A: entra 3, sale $0 + 3 = 3$ \checkmark; \
        B: entra $7 + 0 = 7$, sale $0 + 4 + 3 = 7$ \checkmark; \
        C: entra $4 + 0 = 4$, sale 4 \checkmark; \
        D: entra $3 + 4 + 1 = 8$, sale 8 \checkmark; \
        E: entra $3 + 4 = 7$, sale $1 + 6 = 7$ \checkmark
  \item \textbf{Sale de O} $3 + 7 + 4 = 14$ y \textbf{llega a T} $8 + 6 = 14$ \checkmark
  \item \textbf{Corte mínimo:} desde O se alcanzan \{O, A, B, C, E\}. Cruzan A$\to$D (3), B$\to$D (4),
        E$\to$D (1) y E$\to$T (6): capacidad $3 + 4 + 1 + 6 = 14 =$ flujo máximo \checkmark
\end{itemize}

\subsection*{Interpretación gráfica}
\begin{center}
\begin{tikzpicture}[x=1.8cm, y=1.12cm]
  \nodosSeervada
  \ar[usado]{O}{A}{3/5}
  \ar[usado]{O}{B}{7/7}
  \ar[usado]{O}{C}{4/4}
  \ar[libre]{A}{B}{0/1}
  \ar[usado, lbl/.style={pos=0.35}]{A}{D}{3/3}
  \ar[libre]{B}{C}{0/2}
  \ar[usado, lbl/.style={pos=0.4}]{B}{D}{4/4}
  \ar[usado, lbl/.style={pos=0.6}]{B}{E}{3/5}
  \ar[usado]{C}{E}{4/4}
  \ar[usado]{D}{T}{8/9}
  \ar[usado, lbl/.style={pos=0.25}]{E}{D}{1/1}
  \ar[usado, lbl/.style={pos=0.65}]{E}{T}{6/6}
  \draw[corte] plot[smooth] coordinates {(2.75,0.2) (2.95,-1.2) (3.4,-1.95) (4.1,-2.3) (4.55,-3.7)};
  \node[font=\small, right] at (4.55,-3.7) {corte mínimo $= 14$};
\end{tikzpicture}
\end{center}
La línea discontinua deja de un lado a O, A, B, C y E, y del otro a D y T. Los cuatro caminos que la cruzan van
llenos y entre todos suman 14: ese es el embudo del parque. En cambio, D$\to$T (8 de 9), O$\to$A (3 de 5) y
B$\to$E (3 de 5) tienen espacio de sobra, así que ampliarlos no serviría. Para tener más viajes habría que ampliar
alguno de los cuatro caminos del corte.

% =====================================================================
\section{Red de Taha: cinco nodos}
\textit{Fuente: adaptado de Taha, Investigación de Operaciones, capítulo de modelos de redes (flujo máximo).}

\subsection*{Planteamiento del problema}
Una red de tuberías lleva un líquido desde el nodo \textbf{1} (fuente) hasta el nodo \textbf{5} (sumidero),
pasando por los nodos 2, 3 y 4. Entre los nodos 3 y 4 hay tuberías en los dos sentidos. Las capacidades, en miles
de litros por hora, son:

\begin{center}\small
\begin{tabular}{lccccccccc}
\toprule
\textbf{Tubería} & 1$\to$2 & 1$\to$3 & 1$\to$4 & 2$\to$3 & 2$\to$5 & 3$\to$4 & 3$\to$5 & 4$\to$3 & 4$\to$5 \\
\midrule
\textbf{Capacidad} & 20 & 30 & 10 & 40 & 30 & 10 & 20 & 5 & 20 \\
\bottomrule
\end{tabular}

\medskip
\begin{minipage}{\linewidth}\centering
\begin{tikzpicture}[x=2.1cm, y=1.0cm]
  \nodosTaha
  \ar{t1}{t2}{0/20}
  \ar{t1}{t3}{0/30}
  \ar{t1}{t4}{0/10}
  \ar{t2}{t3}{0/40}
  \ar{t2}{t5}{0/30}
  \ar[arco, bend left=40, lbl/.style={right}]{t3}{t4}{0/10}
  \ar[arco, bend left=40, lbl/.style={left}]{t4}{t3}{0/5}
  \ar{t3}{t5}{0/20}
  \ar{t4}{t5}{0/20}
\end{tikzpicture}

Red inicial: todo el flujo en 0.
\end{minipage}
\end{center}

La pregunta: \textbf{¿cuánto líquido por hora puede llegar como máximo al nodo 5?}

\subsection*{Formulación del modelo}
\[
\begin{aligned}
\text{Max } Z &= f_{12} + f_{13} + f_{14} \nota{lo que sale de la fuente}\\
\text{s.a. } \quad f_{ij} &\le u_{ij} \ \text{(la tabla)}, \qquad f_{ij} \ge 0\\
f_{12} - f_{23} - f_{25} &= 0 \nota{nodo 2}\\
f_{13} + f_{23} + f_{43} - f_{34} - f_{35} &= 0 \nota{nodo 3}\\
f_{14} + f_{34} - f_{43} - f_{45} &= 0 \nota{nodo 4}
\end{aligned}
\]

\subsection*{Caminos de aumento}
Desde 1 se llega en un paso a 2, 3 y 4, y cada uno tiene un arco directo a 5. Así que los tres primeros caminos
tienen solo 2 arcos, y se toman en el orden de la lista.

\paso{Camino 1: \ $1 \to 2 \to 5$}
Es el primero de los tres: por 1$\to$2 caben 20 y por 2$\to$5 caben 30.
\[
\delta = \min(\underbrace{20}_{12},\ \underbrace{30}_{25}) = 20,
\qquad \text{flujo total} = 0 + 20 = 20 \nota{1$\to$2 se llena}
\]

\paso{Camino 2: \ $1 \to 3 \to 5$}
1$\to$2 ya está lleno; desde 1 se pasa a 3, y de 3 se va directo a 5.
\[
\delta = \min(\underbrace{30}_{13},\ \underbrace{20}_{35}) = 20,
\qquad \text{flujo total} = 20 + 20 = 40 \nota{3$\to$5 se llena; a 1$\to$3 le quedan 10}
\]

\paso{Camino 3: \ $1 \to 4 \to 5$}
Ahora 3$\to$5 también está lleno; queda el camino que pasa por 4.
\[
\delta = \min(\underbrace{10}_{14},\ \underbrace{20}_{45}) = 10,
\qquad \text{flujo total} = 40 + 10 = 50 \nota{1$\to$4 se llena; a 4$\to$5 le quedan 10}
\]

\paso{Camino 4: \ $1 \to 3 \to 4 \to 5$}
Desde 1 solo se llega a 3 (1$\to$2 y 1$\to$4 están llenos). Desde 3 no se puede ir directo a 5, pero sí a 4,
y desde 4 a 5. El camino tiene 3 arcos:
\[
\delta = \min(\underbrace{10}_{13},\ \underbrace{10}_{34},\ \underbrace{10}_{45}) = 10,
\qquad \text{flujo total} = 50 + 10 = 60 \nota{los tres tramos se llenan a la vez}
\]

\begin{center}
\begin{tikzpicture}[x=2.1cm, y=1.1cm]
  \nodosTaha
  \ar[usado]{t1}{t2}{20/20}
  \ar[camino]{t1}{t3}{30/30}
  \ar[usado]{t1}{t4}{10/10}
  \ar[libre]{t2}{t3}{0/40}
  \ar[usado]{t2}{t5}{20/30}
  \ar[camino, bend left=40, lbl/.style={right}]{t3}{t4}{10/10}
  \ar[libre, bend left=40, lbl/.style={left}]{t4}{t3}{0/5}
  \ar[usado]{t3}{t5}{20/20}
  \ar[camino, lbl/.style={pos=0.62}]{t4}{t5}{20/20}
\end{tikzpicture}

\small Después del camino 4 ($1\to3\to4\to5$ resaltado). Flujo total 60.
\end{center}

\paso{¿Hay otro camino?}
Los tres arcos que salen de la fuente van llenos, y como ningún arco llega a 1, tampoco hay flujo que se pueda
devolver desde 1. Desde 1 \textbf{no se llega a ningún otro nodo}: el método termina.

\subsection*{Resultado}
\[
\res{\text{Flujo máximo} = 60}
\]
\begin{itemize}[itemsep=2pt]
  \item Al nodo 5 llegan como máximo \textbf{60 mil litros por hora}.
  \item Flujo por tubería: 1$\to$2 = 20, 1$\to$3 = 30, 1$\to$4 = 10, 2$\to$3 = 0 (sin usar), 2$\to$5 = 20,
        3$\to$4 = 10, 3$\to$5 = 20, 4$\to$3 = 0 (sin usar), 4$\to$5 = 20.
\end{itemize}

\subsection*{Comprobación}
\begin{itemize}[itemsep=2pt]
  \item \textbf{Capacidades:} ningún flujo pasa de su capacidad; 2$\to$5 es la única tubería usada que no va
        llena (20 de 30) \checkmark
  \item \textbf{Conservación:} nodo 2: entra 20, sale $0 + 20 = 20$ \checkmark; \
        nodo 3: entra $30 + 0 + 0 = 30$, sale $10 + 20 = 30$ \checkmark; \
        nodo 4: entra $10 + 10 = 20$, sale $0 + 20 = 20$ \checkmark
  \item \textbf{Sale de 1} $20 + 30 + 10 = 60$ y \textbf{llega a 5} $20 + 20 + 20 = 60$ \checkmark
  \item \textbf{Corte mínimo:} desde 1 solo se alcanza $\{1\}$. Cruzan 1$\to$2 (20), 1$\to$3 (30) y 1$\to$4 (10):
        capacidad $20 + 30 + 10 = 60 =$ flujo máximo \checkmark
\end{itemize}

\subsection*{Interpretación gráfica}
\begin{center}
\begin{tikzpicture}[x=2.1cm, y=1.1cm]
  \nodosTaha
  \ar[usado, lbl/.style={pos=0.6}]{t1}{t2}{20/20}
  \ar[usado, lbl/.style={pos=0.6}]{t1}{t3}{30/30}
  \ar[usado, lbl/.style={pos=0.6}]{t1}{t4}{10/10}
  \ar[libre]{t2}{t3}{0/40}
  \ar[usado]{t2}{t5}{20/30}
  \ar[usado, bend left=40, lbl/.style={right}]{t3}{t4}{10/10}
  \ar[libre, bend left=40, lbl/.style={left}]{t4}{t3}{0/5}
  \ar[usado]{t3}{t5}{20/20}
  \ar[usado, lbl/.style={pos=0.62}]{t4}{t5}{20/20}
  \draw[corte] (0.55,0.1) -- (0.55,-4.1) node[below, font=\small] {corte mínimo $= 60$};
\end{tikzpicture}
\end{center}
Aquí el embudo está \textbf{justo a la salida de la fuente}: de 1 solo salen tres tuberías y entre todas aguantan
$20 + 30 + 10 = 60$. Dentro de la red sobra espacio (2$\to$3 ni se usa y 2$\to$5 va a dos tercios), pero eso no
ayuda porque no se puede meter más de lo que sale de 1. Para mandar más habría que ampliar las tuberías que salen
de la fuente.

\bigskip
\section*{Resumen de los tres ejercicios}
\begin{center}\small
\begin{tabular}{lccp{5.0cm}c}
\toprule
\textbf{Ejercicio} & \textbf{Red} & \textbf{Caminos} & \textbf{Corte mínimo} & \textbf{Flujo máximo} \\
\midrule
Sunco Oil & 5 nodos, 6 arcos & 2 & las 2 tuberías que llegan a si & 3 \\
Parque Seervada & 7 nodos, 12 arcos & 5 & A$\to$D, B$\to$D, E$\to$D y E$\to$T & 14 \\
Red de Taha & 5 nodos, 9 arcos & 4 & las 3 tuberías que salen de 1 & 60 \\
\bottomrule
\end{tabular}
\end{center}

\begin{idea}[Conclusiones]
\begin{itemize}[itemsep=2pt, leftmargin=*]
  \item El flujo máximo es lo más que se puede mandar de la fuente al sumidero sin pasarse de ninguna capacidad y
        respetando que en cada nodo intermedio lo que entra, sale.
  \item Ford-Fulkerson repite tres cosas: buscar un camino con espacio, ver su cuello de botella y mandarlo.
        Con la regla de Edmonds y Karp se toma siempre el camino con menos arcos.
  \item La red residual guarda lo que todavía cabe y lo que se puede devolver; las devoluciones corrigen caminos
        mal escogidos.
  \item El método termina cuando ya no se llega al sumidero. Los nodos alcanzables marcan el corte mínimo, y su
        capacidad es igual al flujo máximo: esa es la comprobación.
  \item El corte mínimo dice dónde está el embudo y, por lo tanto, qué tramos vale la pena ampliar.
\end{itemize}
\end{idea}

\textbf{Verificación:} los tres ejercicios se comprobaron con la calculadora (\emph{Simplex Arcade}, modo 4P Flujo
máximo, memorias 10, 11 y 12), que reproduce los mismos caminos de aumento, cuellos de botella y cortes. Además,
cada red se planteó como programación lineal y se resolvió con el Simplex Revisado y con la librería
javascript-lp-solver: los tres dan el mismo flujo máximo (3, 14 y 60).

\section*{Referencias}
\addcontentsline{toc}{part}{Referencias}
\begin{itemize}[leftmargin=*, itemsep=3pt]
  \item Edmonds, J. y Karp, R. M. (1972). Theoretical improvements in algorithmic efficiency for network flow problems. \textit{Journal of the ACM}, 19(2).
  \item Ford, L. R. y Fulkerson, D. R. (1956). Maximal flow through a network. \textit{Canadian Journal of Mathematics}, 8.
  \item Hillier, F. S. y Lieberman, G. J. \textit{Introducción a la Investigación de Operaciones}. McGraw-Hill. Capítulo de optimización de redes.
  \item Taha, H. A. \textit{Investigación de Operaciones}. Pearson. Capítulo de modelos de redes.
  \item Winston, W. L. \textit{Investigación de operaciones: aplicaciones y algoritmos}. Thomson. Capítulo de modelos de redes.
\end{itemize}

\end{document}
